%% ============================================================
%% SECTION: Our Proposal
%% ============================================================

\section{Our Proposal}

In this work, we adopt a Three-Dimensional Convolutional Neural Network (3D-CNN) framework tailored for the segmentation of volumetric brain tumors. Unlike Two-Dimensional slice-based methods, our model operates directly on cropped volumes of size $128\times128\times128$, thus preserving the spatial context across all three dimensions. The network input consists of three MRI modalities: T1 contrast enhanced (T1ce), T2 and FLAIR stacked as channels, while the T1 modality is excluded to avoid redundancy, as T1ce provides most of the anatomical information available in T1 while additionally highlighting contrast-enhanced regions. Our Convolutional approach is built on the classical U-Net shape backbone where there is an encoder (contracting path) and a decoder (expanding path) connected through a bottleneck and skip connections(residual connections). But in our approach we replace the classical $3\times3\times3$ convolution used for the basic U-net with custom Dense Inception Residual Convolution (DRIC) blocks to capture multi-scale features and attention gates to selectively emphasize tumor regions during feature propagation through skip connections. Although computationally more demanding, 3D convolutions were adopted to preserve volumetric context rather than segmenting each 2D slice as in our literature review, which tends to capture more volumetric context. To address the high memory requirements of 3D volumetric images by cropping into the data and discarding the empty spaces present. Now, there are few major components that make up our approach in building this architecture.
